% Toponymia Europaea — Research Paper Template
% Auto-generated by toponymia paper generator
\documentclass[12pt,a4paper]{article}

\usepackage[utf8]{inputenc}
\usepackage[T1]{fontenc}
\usepackage{lmodern}
\usepackage[english]{babel}
\usepackage{amsmath,amssymb}
\usepackage{graphicx}
\usepackage{booktabs}
\usepackage{hyperref}
\usepackage[margin=2.5cm]{geometry}
\usepackage{natbib}
\usepackage{float}

\title{<TITLE>}
\author{<AUTHOR>}
\date{\today}

\begin{document}

\maketitle

\begin{abstract}
<ABSTRACT>
\end{abstract}

\section{Introduction}

This paper presents a quantitative analysis of European place names using the
Toponymia Europaea framework. The analysis is fully reproducible: all data,
code, and statistical tests are available at
\url{https://github.com/egkristi/ToponymiaEuropaea}.

\section{Data}

\begin{table}[H]
\centering
\caption{Dataset summary}
\begin{tabular}{lrr}
\toprule
Country & Records & Sources \\
\midrule
<COUNTRY_TABLE>
\bottomrule
\end{tabular}
\end{table}

\subsection{Data Sources}

<SOURCES_DESCRIPTION>

\section{Methods}

The analysis employs the following statistical framework:
\begin{itemize}
    \item Spatial autocorrelation (Moran's I) for geographic clustering
    \item Permutation tests ($n=10{,}000$) for element distribution significance
    \item Correspondence analysis for language--geography associations
\end{itemize}

All tests use $\alpha = 0.05$ with Bonferroni correction for multiple comparisons.

\section{Results}

<RESULTS>

\section{Discussion}

<DISCUSSION>

\section{Conclusion}

<CONCLUSION>

\section*{Reproducibility Statement}

This analysis was generated using Toponymia Europaea v0.3.0. To reproduce:
\begin{verbatim}
git clone https://github.com/egkristi/ToponymiaEuropaea.git
cd ToponymiaEuropaea
uv sync --extra dev
uv run pytest  # verify installation
uv run toponymia analyze <COMMAND>
\end{verbatim}

\bibliographystyle{apalike}
\bibliography{references}

\end{document}
